% RSF 2026 — Lucrare de cercetare
\documentclass[12pt, a4paper]{article}

\usepackage[T1]{fontenc}
\usepackage[utf8]{inputenc}
\usepackage[romanian]{babel}
\usepackage{lmodern}
\usepackage[margin=2.5cm]{geometry}
\usepackage{setspace}
\usepackage{parskip}
\usepackage{graphicx}
\usepackage{float}
\usepackage{caption}
\usepackage{booktabs}
\usepackage{xcolor}
\usepackage{hyperref}
\usepackage{url}
\usepackage{amsmath, amssymb}
\usepackage{enumitem}
\usepackage[style=apa, backend=biber, sorting=nyt]{biblatex}
\usepackage{csquotes}

\onehalfspacing
\setlength{\parindent}{0pt}
\setlength{\parskip}{6pt}
\hypersetup{colorlinks=true,linkcolor=black,citecolor=blue!60!black,urlcolor=blue!60!black}
\captionsetup{font=small, labelfont=bf}
\addbibresource{references.bib}

\newcommand{\titlulProiectului}{YR Motion Game: controlul unui joc video prin mișcări corporale recunoscute în timp real}
\newcommand{\numeElevi}{Gabriela Benghe, clasa a XI-a A}
\newcommand{\numeMentor}{[de completat]}
\newcommand{\institutie}{Colegiul Național „Ion C. Brătianu”, Pitești}

\begin{document}

\begin{titlepage}
\centering
\vspace*{2cm}
{\Large\textsc{Romanian Science Festival 2026}}\\[0.3cm]
{\small Program de Cercetare pentru Elevi}\\[2.5cm]
{\LARGE\bfseries \titlulProiectului}\\[2cm]
{\large \numeElevi}\\[0.5cm]
{\normalsize Mentor: \numeMentor}\\[0.3cm]
{\normalsize \institutie}\\[3cm]
\vfill
{\normalsize Pitești, septembrie 2026}
\end{titlepage}

\tableofcontents
\newpage

\section{Introducere}

În majoritatea jocurilor video, intenția jucătorului este transmisă prin apăsarea
unui buton. Într-un joc bazat pe mișcare, corpul devine dispozitivul de control, iar
o acțiune fizică produce un rezultat vizibil în lumea virtuală. Legătura este ușor de
urmărit chiar și de cineva care vede jocul pentru prima dată: un braț întins poate
lansa un proiectil, iar o coborâre a corpului poate evita un atac.

Ideea folosirii unei camere obișnuite pentru jocuri corporale există de peste două
decenii. Proiectul \textit{Shadow Boxer} a demonstrat încă din 2004 că mișcarea unei
persoane poate controla un joc prin intermediul unei camere web
\parencite{hoysniemi2004}. Între timp, tehnologia de estimare a posturii a progresat.
BlazePose a fost proiectat pentru urmărirea corpului în timp real pe dispozitive
obișnuite și produce 33 de repere anatomice \parencite{bazarevsky2020}. MediaPipe
Pose Landmarker pune la dispoziție aceleași tipuri de repere în coordonate de imagine
și coordonate spațiale \parencite{googlepose2026}.

Localizarea corpului nu rezolvă automat problema interacțiunii. Oamenii nu execută
două mișcări în exact același mod. Unghiul brațului, poziția cotului, viteza și
distanța față de cameră variază. Un detector prea strict respinge mișcări pe care un
observator le consideră corecte, iar unul prea permisiv poate transforma postura
relaxată într-o comandă. Cercetarea recomandă acceptarea variației firești a corpului
și evitarea preciziei care nu aduce un beneficiu real experienței
\parencite{isbister2015}. O analiză sistematică a literaturii a identificat 107
recomandări pentru proiectarea jocurilor corporale, grupate în teme precum alegerea
mișcărilor, feedbackul, siguranța și accesibilitatea \parencite{subramanian2020}.

YR Motion Game investighează problema într-un caz concret. Este un joc de luptă care
rulează în browser și folosește patru comenzi corporale:

\begin{itemize}[nosep]
    \item \textbf{ATAC}: brațul drept este întins lateral;
    \item \textbf{BLOCARE}: antebrațele sunt ridicate și încrucișate la piept;
    \item \textbf{ESCHIVĂ}: partea superioară a corpului coboară;
    \item \textbf{SPECIALĂ}: ambele brațe sunt ridicate.
\end{itemize}

Sistemul nu solicită înregistrarea unor exemple personale și nu compară utilizatorul
cu o postură perfectă. MediaPipe furnizează reperele anatomice, iar un detector
geometric transparent decide ce comandă descriu acestea.

\section{Obiective și Întrebarea de Cercetare}

\textbf{Obiectivul general} este dezvoltarea unui sistem web capabil să transforme
patru mișcări corporale distincte în acțiuni de joc clare, rapide și independente,
pe cât posibil, de distanța utilizatorului față de cameră.

\textbf{Obiectivele specifice} sunt:

\begin{enumerate}[nosep]
    \item definirea unui detector geometric explicabil, care să accepte variații
    naturale ale aceleiași mișcări;
    \item reducerea dependenței de distanța față de cameră prin normalizarea
    distanțelor cu lățimea umerilor;
    \item realizarea unui flux în care o postură menținută produce o singură comandă,
    iar feedbackul apare numai pentru acțiunile acceptate de joc;
    \item menținerea unei latențe reduse prin evitarea acumulării cadrelor video vechi.
\end{enumerate}

\textbf{Întrebarea de cercetare:} \textit{În ce măsură poate un detector geometric
normalizat la proporțiile corpului să transforme patru posturi aproximative în
comenzi distincte și coerente pentru un joc video în timp real?}

Ipoteza tehnică este că raportarea distanțelor la lățimea umerilor păstrează decizia
detectorului la o redimensionare uniformă a corpului, iar folosirea unei poziții
neutre reduce activările false pentru BLOCARE și ESCHIVĂ.

\section{Metodologie}

\subsection{Tipul cercetării}

Lucrarea reprezintă o cercetare aplicată de tip \textit{design și implementare de
sistem}, însoțită de verificare software. A fost construit un prototip funcțional,
apoi proprietățile sale au fost verificate prin cazuri sintetice și teste deterministe
ale motorului de luptă. Nu a fost realizat încă un experiment controlat cu
participanți; rezultatele descriu comportamentul programului, nu performanța sa
statistică în populația generală.

\subsection{Instrumente de colectare a datelor}

Intrarea principală este fluxul unei camere web solicitat la rezoluția ideală de
$640\times480$ pixeli și aproximativ 30 de cadre pe secundă. MediaPipe Pose
Landmarker identifică reperele corpului. Dintre cele 33 disponibile sunt folosite
numai nasul, umerii, coatele și încheieturile. Urmărirea degetelor nu este necesară.

În modul de dezvoltare, aplicația afișează valorile geometrice, gestul brut,
evenimentul emis și acțiunea acceptată. Sunt urmărite și frecvența camerei, durata
inferenței, cadrele omise și tipul de procesare, GPU sau CPU. Panoul este actualizat
de aproximativ patru ori pe secundă, pentru a nu redesena întreaga interfață la
fiecare cadru.

\subsection{Eșantion / Materiale}

Nu există un eșantion uman în etapa documentată. Materialul de verificare este format
din schelete sintetice cu 33 de puncte, construite pentru a reprezenta posturi
cunoscute. Cazurile pozitive trebuie să producă una dintre cele patru comenzi, iar
cazurile negative includ poziția neutră, brațul drept lăsat în jos și un singur braț
ridicat.

Pentru testarea invarianței la scară, toate punctele unei posturi sunt apropiate de
centrul corpului cu factorul $0{,}6$. Postura rezultată simulează aceeași persoană
aflată la o distanță mai mare de cameră.

\subsection{Normalizarea geometrică}

Notăm cu $U_S$ și $U_D$ umerii stâng și drept. Centrul umerilor $C$ și lățimea
umerilor $s$ sunt:

\begin{align}
C_x &= \frac{U_{Sx}+U_{Dx}}{2}, &
C_y &= \frac{U_{Sy}+U_{Dy}}{2}, \\
s &= \sqrt{(U_{Sx}-U_{Dx})^2+(U_{Sy}-U_{Dy})^2}.
\end{align}

Toate distanțele folosite sunt împărțite la $s$ sau la lățimea umerilor din poziția
neutră, $s_0$. La o scalare uniformă cu factorul $k$, atât distanța măsurată, cât și
lățimea umerilor sunt înmulțite cu $k$, iar raportul rămâne neschimbat. Reperele sunt
acceptate la o vizibilitate de minimum $0{,}20$. Dacă unul dispare foarte scurt,
ultima postură completă poate fi folosită cel mult 160 ms.

\subsection{Poziția neutră}

La început sunt colectate minimum opt cadre valide pe o perioadă de cel puțin 800 ms.
Se calculează medianele pentru centrul și lățimea umerilor, nas, coate și
încheieturi. Referința este apoi înghețată, astfel încât să nu urmeze utilizatorul în
timpul eschivei.

\subsection{Regulile de recunoaștere}

Ordinea de verificare este SPECIALĂ, BLOCARE, ESCHIVĂ, ATAC.

\textbf{SPECIALĂ.} Ridicarea fiecărei încheieturi față de umărul corespunzător,
împărțită la $s$, trebuie să depășească $-0{,}05$. Toleranța permite coate îndoite și
încheieturi foarte puțin sub umeri.

\textbf{BLOCARE.} Încheieturile trebuie să fie în zona superioară a trunchiului,
ambele brațe trebuie ridicate față de neutru, iar antebrațele trebuie orientate spre
centru. Încheieturile trebuie să fie la mai puțin de o lățime de umeri una de alta
sau să indice o încrucișare. Corpul rămâne aproximativ drept.

\textbf{ESCHIVĂ.} Cu $C_{y0}$, $N_{y0}$ și $s_0$ din poziția neutră:

\begin{align}
d_{umeri} &= \frac{C_y-C_{y0}}{s_0}, &
d_{nas} &= \frac{N_y-N_{y0}}{s_0}.
\end{align}

Eschiva este detectată când $d_{umeri}>0{,}25$ sau $d_{nas}>0{,}28$. Poziția
brațelor nu influențează rezultatul.

\textbf{ATAC.} Sunt folosite umărul și încheietura brațului drept anatomic:

\begin{align}
e_x &= \frac{|M_{Dx}-U_{Dx}|}{s}, &
e_y &= \frac{|M_{Dy}-U_{Dy}|}{s}, \\
\alpha &= \operatorname{atan2}(|M_{Dy}-U_{Dy}|,|M_{Dx}-U_{Dx}|).
\end{align}

Atacul este detectat dacă $e_x>0{,}65$, $e_y<0{,}60$ și $\alpha<40^\circ$.
Valoarea absolută face regula independentă de oglindirea imaginii.

\subsection{Emiterea și acceptarea acțiunii}

Gestul trebuie să rămână stabil 100 ms pentru a emite un singur eveniment. Cât timp
postura este menținută, nu se emit altele. Detectorul se reactivează după ce gestul
lipsește 200 ms.

Evenimentul este evaluat de motorul de luptă. Numai o acțiune acceptată poate porni
animația, crea un proiectil, activa apărarea, afișa eticheta și intra în istoricul
reușitelor. Gestul SPECIALĂ executat înainte ca energia să ajungă la 100\% este
recunoscut, dar nu produce atacul special.

\subsection{Controlul latenței}

Pose Landmarker rulează cel mult o dată la 33 ms. Același timp video nu este procesat
de două ori. Dacă o inferență este activă, cadrul următor este omis, nu adăugat într-o
coadă. Aplicația folosește \texttt{requestVideoFrameCallback} atunci când este
disponibil și \texttt{requestAnimationFrame} ca rezervă. Inițializarea încearcă
procesarea GPU și revine la CPU în caz de eșec.

\section{Rezultate și Discuții}

\subsection{Rezultate}

Detectorul de producție a fost verificat cu posturi sintetice. Tabelul
\ref{tab:gesturi} prezintă principalele cazuri confirmate de funcțiile de test.

\begin{table}[H]
    \centering
    \caption{Cazuri de verificare pentru detectorul geometric}
    \label{tab:gesturi}
    \begin{tabular}{p{0.49\textwidth}p{0.23\textwidth}p{0.16\textwidth}}
        \toprule
        \textbf{Caz testat} & \textbf{Rezultat} & \textbf{Verificat} \\
        \midrule
        Ambele brațe ridicate & SPECIALĂ & Da \\
        Antebrațe în X la piept & BLOCARE & Da \\
        Corp coborât față de neutru & ESCHIVĂ & Da \\
        Braț drept întins lateral & ATAC & Da \\
        Același atac redimensionat la 60\% & ATAC & Da \\
        Același X redimensionat la 60\% & BLOCARE & Da \\
        Braț drept lateral în sens opus pe axa $x$ & ATAC & Da \\
        Poziții neutre simetrice și asimetrice & Nicio comandă & Da \\
        Braț drept în jos sau vertical & Nu ATAC & Da \\
        \bottomrule
    \end{tabular}
\end{table}

Testele temporale confirmă că o postură menținută produce un singur eveniment, iar
revenirea la neutru urmată de repetare produce un eveniment nou. Motorul de luptă a
fost verificat separat. Fiecare proiectil poate produce cel mult un impact. La
contactul unui proiectil advers, ordinea este ESCHIVĂ, BLOCARE, apoi daună normală.
O eschivă activă produce zero daune. O blocare activă creează un proiectil reflectat,
iar adversarul primește daune numai la sosirea acestuia. O apărare expirată nu oferă
imunitate.

\begin{table}[H]
    \centering
    \caption{Configurația principală a luptei}
    \label{tab:config}
    \begin{tabular}{lc}
        \toprule
        \textbf{Parametru} & \textbf{Valoare} \\
        \midrule
        Viața jucătorului & 100 \\
        Viața adversarilor & 90 / 115 / 105 / 135 \\
        Daună ATAC & 15 \\
        Daună SPECIALĂ & 30 \\
        Daună reflectată & 12 \\
        Dauna adversarilor & 10 / 12 / 14 / 16 \\
        Interval atac advers & 900--2400 ms \\
        Probabilitate eschivă adversar & 0,25 \\
        Pauză între eschivele adversarului & 1800 ms \\
        \bottomrule
    \end{tabular}
\end{table}

În etapele avansate, adversarul estimează următoarea comandă folosind frecvența
acțiunilor acceptate și tranzițiile observate. Mecanismul folosește numai istoricul,
nu gestul aflat în curs. Ideea este compatibilă cu cercetarea despre adaptarea
jocurilor pe baza modelelor de jucător \parencite{yannakakis2009,cowley2016}, dar
implementarea proiectului este intenționat simplă și explicabilă.

\subsection{Discuții}

Rezultatele susțin ipoteza tehnică într-un sens limitat: pentru cazurile sintetice,
redimensionarea uniformă nu schimbă decizia. Explicația este dată de raportarea la
lățimea umerilor, mai robustă decât o distanță brută în pixeli.

Poziția neutră rezolvă o ambiguitate importantă. Apropierea mâinilor de trunchi nu
este suficientă pentru identificarea unei blocări, deoarece poate apărea și în repaus.
Compararea ambelor brațe cu referința neutră și verificarea orientării antebrațelor
fac diferența dintre o gardă activă și o postură relaxată.

Arhitectura separă gestul brut, evenimentul emis și acțiunea acceptată. Separarea
previne un feedback înșelător. Dacă energia specială nu este completă, gestul poate
fi corect, dar personajul nu execută atacul și eticheta de succes nu apare.

Abordarea are limite. Scalarea uniformă nu reproduce rotația corpului, scurtarea
aparentă a unui braț orientat spre cameră, ocluzia, lumina slabă sau îmbrăcămintea
largă. Testele sintetice demonstrează proprietăți ale formulelor, nu o rată de
recunoaștere pentru populația generală. Nici performanța nu a fost încă măsurată
controlat pe mai multe dispozitive.

O evaluare viitoare ar trebui să includă persoane cu proporții corporale și
experiențe de joc diferite. Fiecare gest ar fi repetat la două distanțe marcate.
Rezultatele ar trebui raportate separat prin precizie, sensibilitate, scor F1 și
matrice de confuzie. Latența ar trebui descrisă prin mediană și percentila 95, deoarece
o medie poate ascunde întârzierile rare, dar deranjante. Utilizabilitatea poate fi
măsurată separat prin observație și System Usability Scale \parencite{brooke1996}.

\section{Concluzii}

Întrebarea de cercetare a primit un răspuns tehnic preliminar: detectorul geometric
normalizat separă cele patru posturi în cazurile de test și își păstrează decizia la
redimensionarea uniformă a coordonatelor. Rezultatul nu dovedește aceeași performanță
pentru toate persoanele și toate condițiile de filmare.

Contribuția principală este integrarea coerentă a reperelor MediaPipe, normalizării
prin lățimea umerilor, poziției neutre, filtrării temporale și acceptării acțiunii de
către motor. Împreună, aceste mecanisme transformă postura într-o comandă stabilă și
într-o reacție vizuală care corespunde logicii luptei.

Direcția imediată de cercetare este evaluarea cu utilizatori reali. Aceasta trebuie să
măsoare separat fiecare gest, activările false în poziție neutră, efectul distanței,
latența și claritatea interacțiunii. Pragurile pot fi apoi ajustate pe baza datelor,
fără a pierde transparența detectorului.

\newpage
\printbibliography[title={Referințe Bibliografice}]

\end{document}
